\documentclass[12pt,a4paper]{ctexart}
\usepackage[top=2.5cm,bottom=2.5cm,left=2.5cm,right=2.5cm]{geometry}
\usepackage{amsmath,amssymb,amsthm,physics}
\usepackage{graphicx,float,booktabs,array,caption}
\usepackage{hyperref,xcolor,enumitem}
\usepackage[numbers,sort&compress]{natbib}
\hypersetup{colorlinks=true,linkcolor=blue,citecolor=blue,urlcolor=blue}
\theoremstyle{definition}
\newtheorem{definition}{定义}[section]
\newtheorem{remark}{注记}[section]

\title{\heiti 格点QCD中的重夸克有效理论（HQET）：\\
从重夸克对称性到粲物理的连续极限计算}
\author{基于本库文献的综合解析}
\date{\today}

\begin{document}
\maketitle

\begin{abstract}
\noindent
重夸克有效理论（Heavy Quark Effective Theory, HQET）是现代QCD中处理含重夸克（粲$c$、底$b$）强子系统的核心有效场论框架。当重夸克质量$m_Q \gg \Lambda_{\text{QCD}}$时，重夸克的自旋和味自由度在领头阶退耦——即\textbf{重夸克对称性}（heavy quark symmetry）——使得含重夸克的强子物理具有普适的简化结构。HQET在格点QCD中有三重角色：（1）作为处理$O(m_Q a)$离散化误差的理论框架——通过重夸克场的重定义系统吸收这些修正；（2）作为LaMET的形式类比原型——季向东在2014年建立LaMET时明确以HQET为蓝本；（3）作为重味物理中因子化定理的基础——通过``两步匹配''（LaMET$+$boosted HQET）从格点准分布振幅提取HQET光锥分布振幅。本文从HQET的基本拉氏量和重夸克对称性出发，详细阐述其在格点QCD中的三重应用：以CLQCD合作组在十三个系综上对粲介子质量和衰变常数的连续极限计算为例，展示重夸克改进的正规化和$O(m^2_Q a^2)$离散化误差的系统控制；以LPC合作组的重介子光锥分布振幅计算为例，详述LaMET $\to$ QCD LCDA $\to$ boosted HQET LCDA的``两步匹配''方案；最后讨论HQET作为LaMET理论模板的历史角色及其在$B$物理和味物理中的广阔应用前景。
\end{abstract}

\tableofcontents
\newpage

\section{引言：从重夸克的解耦到HQET}

在QCD中，当夸克质量$m_Q$远大于QCD的固有标度$\Lambda_{\text{QCD}} \sim 300$~MeV时，出现了一个重要的物理简化：重夸克的反粒子产生被指数压低（$\sim e^{-m_Q/\Lambda_{\text{QCD}}}$），且重夸克的自旋和味自由度在$m_Q \to \infty$极限下\textbf{退耦}。这一现象被系统地表述为\textbf{重夸克有效理论}（HQET）。

在标准模型中，粲夸克（$m_c \approx 1.3$~GeV）和底夸克（$m_b \approx 4.2$~GeV）都满足$m_Q \gg \Lambda_{\text{QCD}}$的条件。因此，HQET是理解$D$介子、$B$介子、粲偶素（charmonium）和底偶素（bottomonium）物理的基础理论框架。

\textbf{在格点QCD中，HQET扮演着三重关键角色}，这正是本文的论述主线：

\begin{enumerate}
    \item \textbf{重夸克离散化误差的系统控制：}在格点间距$a$有限时，重夸克质量$m_Q$引入额外的$O(m_Q a)$和$O(m^2_Q a^2)$离散化误差。HQET提供了一种通过重夸克场的重定义来系统吸收这些修正的理论框架——这是CLQCD合作组在物理$\pi$质量、多格点间距系综上实现粲物理百分级精度的关键。

    \item \textbf{LaMET的形式类比原型：}季向东于2014年明确将LaMET建模为HQET的``动量版''类比——$P^z$之于LaMET，如同$m_Q$之于HQET。这一类比不仅具有启发价值，更深刻地揭示了大动量展开作为有效场论的严格数学结构。

    \item \textbf{重味因子化的基础：}在LPC合作组最近的工作中，HQET被用作``两步匹配''方案的第二阶段——先从格点准分布振幅通过LaMET匹配到QCD光锥分布振幅（LCDA），再从QCD LCDA通过boosted HQET匹配到HQET LCDA。
\end{enumerate}

\section{HQET的基本理论结构}

\subsection{重夸克极限与重夸克对称性}

HQET的拉氏量通过对QCD拉氏量在$m_Q \to \infty$极限下的系统展开得到。将重夸克场$Q(x)$分解为``大分量''$h_v(x)$和``小分量''$H_v(x)$：

\begin{equation}
Q(x) = e^{-im_Q v\cdot x} [h_v(x) + H_v(x)],
\label{eq:heavy_quark_decomp}
\end{equation}

其中$v^\mu$是重强子的四速度（$v^2 = 1$），$e^{-im_Q v\cdot x}$因子提取了重夸克质量引起的快速振荡。$h_v$满足投影条件$\slashed{v} h_v = h_v$，$H_v$满足$\slashed{v} H_v = -H_v$。

在$m_Q \to \infty$极限下，$H_v$分量被积掉，得到的HQET有效拉氏量为\cite{Ji2014LaMET}：

\begin{equation}
\mathcal{L}_{\text{HQET}} = \bar{h}_v (iv \cdot D) h_v + \frac{1}{2m_Q} \bar{h}_v (iD_\perp)^2 h_v + \frac{g}{4m_Q} \bar{h}_v \sigma_{\mu\nu} G^{\mu\nu} h_v + \mathcal{O}\left(\frac{1}{m_Q^2}\right)。
\label{eq:HQET_lagrangian}
\end{equation}

\textbf{领头阶的物理含义：}
\begin{itemize}
    \item \textbf{第0阶（$m_Q^0$）：}重夸克的传播子仅依赖于其四速度$v^\mu$——\textbf{自旋和味完全退耦}。所有含单个重夸克的强子（如$D$和$B$介子）在$m_Q \to \infty$极限下具有完全相同的动力学。这就是\textbf{重夸克自旋-味对称性}（heavy quark spin-flavor symmetry）。

    \item \textbf{第1阶（$1/m_Q$）：}\textbf{动能项}$(iD_\perp)^2/(2m_Q)$描述了重夸克在胶子场中的残余动能；\textbf{色磁项}$g\sigma_{\mu\nu}G^{\mu\nu}/(4m_Q)$描述了重夸克自旋与胶子场的相互作用——这是$D$和$D^*$（或$B$和$B^*$）介子质量分裂的来源。

    \item \textbf{更高阶：}系统的$1/m_Q$展开提供了处理重味物理中幂次修正的理论框架。
\end{itemize}

\subsection{HQET的标度分离与因子化}

HQET作为一种有效的因子化工具，通过系统地分离标度$m_Q$和$\Lambda_{\text{QCD}}$来实现物理预言的简化。一个一般的物理观测量$\mathcal{O}$可以展开为\cite{Ji2014LaMET}：

\begin{definition}[HQET展开]\cite{Ji2014LaMET}
\begin{equation}
\boxed{\mathcal{O}(m_Q/\Lambda) = Z(m_Q/\Lambda, \Lambda/\mu) \, o(\mu) + \mathcal{O}\left(\frac{1}{m_Q}\right) + \ldots},
\label{eq:HQET_expansion}
\end{equation}
其中$\mathcal{O}$是全理论（QCD）中的观测量，$o(\mu)$是有效理论（HQET，$m_Q \to \infty$）中的对应量，$Z$是微扰可计算的匹配系数。
\end{definition}

\textbf{匹配系数$Z$的物理意义：}$Z$吸收了全理论和有效理论之间所有不同的紫外物理——包括对数发散（$\ln m_Q$）和常数项。$Z$是红外安全的，因此\textbf{可以在微扰QCD中可靠地计算}。非微扰物理完全编码在$o(\mu)$中，且$o(\mu)$是\textbf{味和自旋普适的}。

\section{HQET在格点离散化误差控制中的应用}

\subsection{重夸克质量带来的$O(m_Q a)$问题}

在格点QCD中，离散化误差通常以$(a\Lambda_{\text{QCD}})^2$的形式出现。然而，当存在重夸克时，出现了更严重的修正：$O(m_Q a)$和$O(m^2_Q a^2)$。对于粲夸克（$m_c \approx 1.3$~GeV）在$a \approx 0.1$~fm $\approx 1/(2\text{ GeV})$的格点上：

\begin{equation}
m_c a \approx \frac{1.3\text{ GeV}}{2\text{ GeV}} \approx 0.65, \qquad (m_c a)^2 \approx 0.42。
\end{equation}

如果不对这些$m_Q a$效应进行特殊处理，即使$a \to 0$的连续极限外推也会因收敛缓慢而失去精度\cite{CLQCD2024charm}。

\subsection{CLQCD合作组的重夸克改进方案}

CLQCD合作组\cite{CLQCD2024charm,CLQCD2025quark_mass}提出了一种系统性的重夸克改进方法，其核心思想类似于HQET：通过\textbf{重夸克场的重定义}来吸收有限$a$处的$m_Q a$依赖性。

\textbf{方法一：重夸克改进的正规化}\cite{CLQCD2024charm}

在连续极限下，矢量流守恒条件$\langle H|V_4|H\rangle / \langle H|H\rangle = 1$精确成立，且$Z_V$应与强子态$H$的选择无关。但在有限格点间距下，$Z_V(H)$表现出\textbf{显著的重夸克质量依赖性}。这一依赖性的主要来源是格点上重夸克传播子的$a$依赖性。

解决方案是使用\textbf{重夸克改进的正规化因子}：

\begin{equation}
Z^c_V \equiv Z_V(\eta_c), \qquad Z^c_X \equiv Z^c_V \frac{Z_X}{Z_V},
\label{eq:heavy_improved_Z}
\end{equation}

其中$Z_X/Z_V$的比值在RI/MOM中间方案下计算，并通过手征外推至手征极限以消除轻夸克质量的残余效应。对于味改变流$Y \equiv \bar{c}\Gamma l$，改进的正规化因子推广为$Z^c_Y \equiv \sqrt{Z_V Z^c_V} \cdot Z_Y/Z_V$。

\textbf{方法二：重夸克质量的吸收}\cite{CLQCD2024charm}

在HQET精神的指引下，粲夸克质量在有限格点间距下的$O(m^2_c a^2)$和$O(m^4_c a^4)$修正可以通过重定义粲夸克质量参数来吸收——类似于手征微扰论中将$F_\pi(0)$替换为$F_\pi(m_\pi)$以改善展开收敛性的做法。

\subsection{数值验证：十三个系综上的粲物理}

CLQCD合作组\cite{CLQCD2024charm}在\textbf{十三个$N_f=2+1$味系综}上验证了这些改进方案的有效性，系综参数覆盖：

\begin{table}[H]
\centering
\caption{CLQCD粲介子计算系综参数\cite{CLQCD2024charm}}
\label{tab:charm_ensembles}
\begin{tabular}{cccc}
\toprule
参数 & 范围 \\
\midrule
格点间距 $a$ & $0.052\text{--}0.105$~fm（5个值） \\
$\pi$介子质量 & $136\text{--}351$~MeV（8个值） \\
格点体积 & $24^3\times 64$ 至 $48^3\times 144$ \\
粲夸克质量参数 & $\tilde{m}^v_c \in [0.055, 0.421]$ \\
\bottomrule
\end{tabular}
\end{table}

\textbf{主要结果：}

\begin{enumerate}
    \item \textbf{质量：}在改进后，$D$, $D^*$, $D_s$, $D^*_s$, $\eta_c$, $J/\psi$, $\chi_{c0}$, $\chi_{c1}$, $h_c$的质量在$a \sim 0.1$~fm处的离散化误差被压制到\textbf{仅几个百分点甚至更低}。

    \item \textbf{衰变常数：}所有$S$波开粲介子（$D$, $D^*$, $D_s$, $D^*_s$）的衰变常数经过重夸克改进后，格点间距依赖性几乎消失，与HPQCD合作组使用HISQ作用量的结果在$1$--$2\%$水平一致。

    \item \textbf{粲夸克质量：}$\overline{\text{MS}}$方案下$m_c(m_c) = 1.2933(72)(95)$~GeV，连续极限外推后的不确定度仅为$\sim 1\%$。

    \item \textbf{超精细分裂：}$D^*_s - D_s$和$J/\psi - \eta_c$的超精细分裂在连续极限下与实验值在$1\sigma$内一致，验证了重夸克改进方案对自旋相关观测量的有效性。

    \item \textbf{P波粲偶素分裂：}$\Delta_{\text{1P-1S}} = 462(14)(09)$~MeV，自旋-轨道分裂$40(18)(06)$~MeV，张量分裂$15.4(7.5)(0.2)$~MeV，均与实验值一致。
\end{enumerate}

\subsection{联合外推方案}

连续极限和物理质量的提取通过以下联合拟合实现\cite{CLQCD2024charm}：

\begin{equation}
X(m_\pi, m_{\eta_s}, a) = X(m^{\text{phys}}_\pi, m^{\text{phys}}_{\eta_s}, 0) + d^X_1 (m^2_\pi - (m^{\text{phys}}_\pi)^2) + d^X_2 (m^2_{\eta_s} - (m^{\text{phys}}_{\eta_s})^2) + d^X_3 a^2 + d^X_4 a^4。
\label{eq:joint_fit}
\end{equation}

其中$a^4$项对于某些观测量（特别是$\overline{\text{MS}}$方案下的粲夸克质量）是\textbf{必需的}——它捕捉了由$O(m^4_c a^4)$效应引起的剩余离散化误差。拟合的$\chi^2/\text{d.o.f} \sim 1$验证了该参数化的充分性。

系统误差通过以下方式估计：
\begin{itemize}
    \item \textbf{拟合形式：}比较加性模型（Eq.\ref{eq:joint_fit}）和乘性模型，差异远小于统计误差；
    \item \textbf{有限体积效应：}在拟合中增补$e^{-m_\pi L}$项，影响可忽略；
    \item \textbf{$O(a\alpha_s)$项：}增补后影响很小；
    \item \textbf{标度参数$w_0$：}涨落FLAG平均值$w_0 = 0.1736(9)$~fm。
\end{itemize}

\section{HQET作为LaMET的理论模板}

\subsection{季向东的类比}

季向东在2014年系统表述LaMET时\cite{Ji2014LaMET}，明确地以HQET为蓝本构建了整个理论框架：

\begin{quote}
``The situation is very similar to that of the heavy quark effective field theory (HQET), which by now has been well-understood and widely used. Here I discuss in some detail this new large momentum effective field theory (LaMET), which allows parton physics to be calculated in lattice QCD at finite $P^z$ in a systematic approximation.''
\end{quote}

\begin{table}[H]
\centering
\caption{HQET与LaMET的精确形式类比\cite{Ji2014LaMET}}
\label{tab:HQET_LaMET_analogy}
\begin{tabular}{lccc}
\toprule
要素 & HQET & LaMET \\
\midrule
\textbf{大标度}          & 重夸克质量 $m_Q$                    & 强子动量 $P^z$ \\
\textbf{有效理论极限}    & $m_Q \to \infty$                     & $P^z \to \infty$ \\
\textbf{有效自由度}      & 重夸克$h_v$ + 轻自由度              & 部分子（光锥关联） \\
\textbf{小参量展开}      & $1/m_Q$（幂次修正）                 & $1/P^z$（高阶twist修正） \\
\textbf{匹配因子}        & $Z(m_Q/\Lambda, \Lambda/\mu)$        & $Z(P^z/\Lambda, \Lambda/\mu)$ \\
\textbf{有效理论观测量}  & $o(\mu)$（HQET矩阵元）              & $f(\mu)$（光锥PDF） \\
\textbf{全理论观测量}    & $\mathcal{O}(m_Q/\Lambda)$（QCD）    & $F(P^z/\Lambda)$（准PDF） \\
\textbf{对称性}          & 重夸克自旋-味对称性                & Lorentz boost不变性（$P^z \to \infty$） \\
\textbf{红外安全性}      & $Z$是红外安全的，非微扰物理在$o$中 & $Z$是红外安全的，非微扰物理在$f$中 \\
\midrule
\textbf{相似性} & \multicolumn{2}{c}{大标度极限下的系统可展开性、微扰匹配、幂次修正} \\
\textbf{差异}   & \multicolumn{2}{c}{HQET：拉氏量层次的EFT；LaMET：算符/矩阵元层次的EFT} \\
\bottomrule
\end{tabular}
\end{table}

\subsection{极限非对易性的共同特征}

HQET和LaMET共享一个深刻的\textbf{场论特征}：极限的非对易性\cite{Ji2014LaMET,Ji2017More}。

\begin{itemize}
    \item \textbf{HQET中：}全理论QCD是先进行UV重整化、后在$m_Q \to \infty$下匹配到HQET。在HQET中，$m_Q \to \infty$的限制在Feynman积分内被取（$\int d^4k / (k^2 - m^2_Q) \to 1/(v\cdot k)$），从而改变了UV行为。
    \item \textbf{LaMET中：}全理论格点QCD是先施加格点UV截断（$a$有限）、后在有限$P^z$下计算。光锥PDF是先取$P^z \to \infty$、后进行UV重整化。两种极限顺序的差异导致了UV行为的不同。
\end{itemize}

在这两种情形中，\textbf{红外物理是相同的}——这是有效场论成立的充要条件。匹配因子$Z$中包含了所有与不同极限顺序相关的UV物理。

\subsection{LaMET中的辅助``重夸克''场方法}

在HQET的启发下，Chen、Ji和Zhang\cite{Chen2016improved}将辅助``重夸克''场方法应用于准PDF中Wilson线的重整化。定义辅助$z$-场$Z(z) = L(z, \infty)$，满足：

\begin{equation}
[\partial_z - igA^z(z)] Z(z) = 0,
\label{eq:Z_field_eom}
\end{equation}

这是\textbf{无限重夸克运动方程的精确类比}——$z$方向扮演了``时间''的角色，$A^z$类似于HQET中的``剩余质量''耦合。Wilson线的重整化：

\begin{equation}
L_{\text{ren}}(z, 0) = Z^{-1}_Z e^{-\delta m |z|} L(z, 0),
\end{equation}

\textbf{完全平行于HQET中重-轻夸克流} $J = \bar{q}\Gamma h_v$ 的重整化。

Zhang等人\cite{Zhang2019}进一步将此方法推广到\textbf{伴随表示}，用于处理胶子准PDF算符的重整化。辅助伴随重夸克场$Q$的拉氏量$\mathcal{L} = \mathcal{L}_{\text{QCD}} + \bar{Q}(in\cdot D - m)Q$具有与HQET完全相同的结构，只是基本表示被替换为伴随表示。

\section{``两步匹配'': HQET在重介子LCDAs计算中的应用}

\subsection{重介子光锥分布振幅的物理重要性}

重介子（$B$, $D$）的光锥分布振幅（LCDAs）是描述重介子内部夸克动量分布的基本非微扰函数。它们在$B$介子非轻子衰变的QCD因子化中起着核心作用\cite{Han2025heavy}——例如$B^0 \to \pi^+\pi^-$的衰变振幅可以因子化为硬散射核与$B$介子和$\pi$介子LCDAs的卷积。

然而，重介子LCDAs的格点计算面临三个根本性困难：

\begin{enumerate}
    \item \textbf{大动量：}LaMET框架要求$P^z \gg \Lambda_{\text{QCD}}$以压制高阶twist修正；
    \item \textbf{重夸克质量：}$m_H \approx m_Q \gg \Lambda_{\text{QCD}}$引入了额外的标度，需要与$P^z$协同处理；
    \item \textbf{光锥定义：}HQET的LCDAs在定义上就涉及光锥关联和HQET有效场，无法直接在格点上模拟。
\end{enumerate}

\subsection{LPC合作组的``两步匹配''方案}

LPC合作组（Han等人\cite{Han2025heavy}）提出了一个巧妙的\textbf{序贯有效理论方法}来克服这些困难。当三个标度满足层级关系：

\begin{equation}
\boxed{P^z \gg m_H \gg \Lambda_{\text{QCD}}},
\label{eq:hierarchy}
\end{equation}

两个大标度$P^z$和$m_H$都落在微扰区域，可以\textbf{分步积掉}：

\begin{center}
\begin{tabular}{cccc}
准分布振幅（格点可计算）& $\xrightarrow{\text{Step I: LaMET（积掉 }P^z\text{）}}$ & QCD LCDA & \\
& & $\xrightarrow{\text{Step II: bHQET（积掉 }m_H\text{）}}$ & HQET LCDA \\
\end{tabular}
\end{center}

\textbf{第一步（LaMET匹配）：}

从格点上计算的准分布振幅$\tilde{\phi}(x, P^z)$出发，通过LaMET因子化匹配到QCD光锥分布振幅$\phi(y, \mu)$：

\begin{equation}
\tilde{\phi}(x, P^z) = \int_0^1 dy\, C\left(x, y, \frac{\mu}{P^z}\right) \phi(y, \mu) + \mathcal{O}\left(\frac{m^2_H}{(P^z)^2}, \frac{\Lambda^2_{\text{QCD}}}{(xP^z, \bar{x}P^z)^2}\right)。
\label{eq:step1_LaMET}
\end{equation}

LPC合作组发现，在领头幂次下，重介子的LaMET匹配核$C^{(1)}_B$与轻介子的完全一致——这是\textbf{重夸克对称性}在起作用的另一个例证。混合-ratio方案中的抵消项为：

\begin{equation}
C^{(1)}_{CT} = -\frac{3\alpha_s C_F}{4\pi} \left[\frac{2\,\text{Si}((x-y)z_s P^z)}{\pi(x-y)}\right]_+。
\label{eq:counter_term}
\end{equation}

\textbf{第二步（boosted HQET匹配）：}

QCD LCDA包含标度远低于$m_H$的软-共线模式（off-shellness $\sim m^2_Q$）。这些模式通过\textbf{boosted HQET}（bHQET）的喷注函数$J$被积掉：

\begin{equation}
\phi(y, \mu) = \int d\omega\, J(y, \omega, \mu, m_H) \, \varphi_+(\omega, \mu) + \mathcal{O}\left(\frac{\Lambda_{\text{QCD}}}{m_H}\right),
\label{eq:step2_bHQET}
\end{equation}

其中$\varphi_+(\omega, \mu)$就是所需的\textbf{HQET LCDA}，$\omega$是轻夸克在重介子静止系中的能量。

在$\omega$大的``尾部''区域（$\omega \sim m_H \gg \Lambda_{\text{QCD}}$），HQET LCDA是\textbf{微扰可计算的}\cite{Han2025heavy}：

\begin{equation}
\varphi_+^{\text{tail}}(\omega, \mu) = \frac{\alpha_s C_F}{\pi\omega} \left[ \left(\frac{1}{2} - \ln\frac{\omega}{\mu}\right) + \frac{4\bar{\Lambda}}{3\omega}\left(2 - \ln\frac{\omega}{\mu}\right) \right]。
\label{eq:tail_region}
\end{equation}

其中$\bar{\Lambda} = m_H - m_Q$刻画了幂次修正的大小。在$\omega$小的``峰区''，非微扰效应主导，格点QCD的输入是不可替代的。

\subsection{数值实现}

LPC合作组在CLQCD的H48P32系综（$a = 0.05187$~fm，体积$48^3 \times 144$，$m_\pi \approx 317$~MeV，549个组态）上进行了验证性计算。粲$D$介子的动量取$P^z = \{0, 2.99, 3.49, 3.98\}$~GeV，使用Coulomb规范固定下的网格源技术进行8784次测量。

经过混合方案重整化和$\lambda = z\cdot P^z$外推后，``两步匹配''得到的HQET LCDAs与现有的唯象模型在形状上一致。模型无关的参数化被提供以便于唯象学应用。

\section{HQET在$B$物理和味物理中的广阔前景}

虽然本库文献主要聚焦于粲物理和LaMET部分子物理中的HQET应用，但HQET最广阔的应用领域在于\textbf{$B$物理和味物理}：

\begin{itemize}
    \item \textbf{$B$介子半轻衰变形状因子：}HQET中的重夸克对称性将所有$B \to D^{(*)}$形状因子在零反冲点（$v \cdot v' = 1$）约化为一个普适的Isgur-Wise函数$\xi(v\cdot v')$。格点QCD对Isgur-Wise函数的精确计算为$|V_{cb}|$的确定提供了关键输入。

    \item \textbf{$B$介子混合与CP破坏：}$B^0$--$\bar{B}^0$混合的袋参数$B_{B}$在HQET框架下具有简化的标度演化行为。格点QCD对$B_{B}$的精确计算是约束CKM幺正三角形的关键输入。

    \item \textbf{重夸克对称性破缺的系统研究：}$1/m_Q$修正（动能项和色磁项）的格点QCD计算为理解$B$和$D$介子谱学和衰变中的精细结构提供了理论基础。

    \item \textbf{HQET在格点上的直接模拟：}在某些方案中，HQET的有效作用量可以直接在格点上离散化——这避免了$O(m_Q a)$离散化误差的困难，但代价是引入了$1/m_Q$展开的截断误差。两种方案的交叉验证是未来的重要方向。
\end{itemize}

\section{总结与展望}

重夸克有效理论（HQET）在格点QCD中发挥着多层面的关键作用。本文的核心结论如下：

\begin{enumerate}
    \item \textbf{重夸克离散化误差的系统控制：}HQET提供了通过重夸克场重定义和重夸克改进正规化来系统吸收$O(m_Q a)$和$O(m^2_Q a^2)$离散化误差的理论框架。CLQCD合作组在十三个系综上的实践\cite{CLQCD2024charm}证明：经过HQET启发的改进后，粲介子质量和$S$波衰变常数在$a \sim 0.1$~fm处的离散化误差被压制到仅几个百分点，粲夸克质量的连续极限精度达到$\sim 1\%$。

    \item \textbf{LaMET的理论模板：}季向东在构建LaMET时明确以HQET为蓝本\cite{Ji2014LaMET}。两者共享相同的有效场论数学结构：大标度极限下的系统展开、微扰可计算的匹配因子、以及幂次压低的修正。这一类比不仅具有启发价值，更赋予了LaMET以严格的有效场论合法性。

    \item \textbf{``两步匹配''方案的物理引擎：}LPC合作组\cite{Han2025heavy}利用HQET（具体为boosted HQET）作为从格点准分布振幅到HQET LCDA的序贯匹配中的第二阶段。这一方案成功解决了长期以来阻碍重介子LCDAs格点计算的三个根本困难（大动量、重夸克质量、光锥定义）。

    \item \textbf{粲物理百分级精度的实现：}CLQCD合作组\cite{CLQCD2024charm,CLQCD2025quark_mass}基于HQET启发的改进方案，在物理$\pi$介子质量、多格点间距的广泛系综上，首次实现了粲介子谱学和衰变常数的百分级精度连续极限计算。这标志着格点重味物理从``可行性验证''到``精度物理''的范式转变。
\end{enumerate}

HQET作为一个成熟的、经过数十年验证的有效场论框架，在格点QCD计算精度不断提升的今天，正在展现出越来越丰富的应用潜力。从粲物理的连续极限计算到$B$物理的形状因子，从LaMET的理论模板到重介子LCDAs的序贯匹配，HQET与格点QCD的深度协同正在推动重味物理进入一个精准预言的新时代。

\begin{thebibliography}{99}

\bibitem{Ji2014LaMET}
X.~Ji,
``Parton physics from large-momentum effective field theory,''
Sci.\ China Phys.\ Mech.\ Astron.\ \textbf{57}, 1407 (2014), arXiv:1404.6680.
\quad\textbf{[来源:} \texttt{docs/Parton Physics from Large-Momentum Effective Field Theory.pdf}\textbf{]}
（含HQET与LaMET的系统类比）

\bibitem{Ji2017More}
X.~Ji, J.-H.~Zhang, and Y.~Zhao,
``More on large-momentum effective theory approach to parton physics,''
Nucl.\ Phys.\ B \textbf{924}, 366 (2017), arXiv:1706.07416.
\quad\textbf{[来源:} \texttt{docs/More On Large-Momentum Effective Theory Approach to Parton Physics.pdf}\textbf{]}

\bibitem{Han2025heavy}
X.-Y.~Han \textit{et al.} (Lattice Parton Collaboration),
``Calculation of heavy meson light-cone distribution amplitudes from lattice QCD,''
(2025), arXiv:2410.18654.
\quad\textbf{[来源:} \texttt{docs/Calculation of heavy meson light-cone distribution amplitudes from lattice QCD.pdf}\textbf{]}
（含LaMET $\to$ QCD LCDA $\to$ boosted HQET LCDA的``两步匹配''方案）

\bibitem{CLQCD2024charm}
H.-Y.~Du \textit{et al.} (CLQCD),
``Charmed meson masses and decay constants in the continuum from the tadpole improved clover ensembles,''
Phys.\ Rev.\ D \textbf{109}, 054507 (2024), arXiv:2310.00814.
\quad\textbf{[来源:} \texttt{docs/Charmed meson masses and decay constants in the continuum from the tadpole improved clover ensembles.pdf}\textbf{]}
（含重夸克改进正规化和$O(m^2_Q a^2)$误差控制）

\bibitem{CLQCD2025quark_mass}
H.-Y.~Du \textit{et al.} (CLQCD),
``Quark masses and low energy constants in the continuum from the tadpole improved clover ensembles,''
Phys.\ Rev.\ D \textbf{111}, 054504 (2025), arXiv:2408.03548.
\quad\textbf{[来源:} \texttt{docs/Quark masses and low energy constants in the continuum from the tadpole improved clover ensembles.pdf}\textbf{]}
（含CLQCD系综参数和粲夸克质量）

\bibitem{Chen2016improved}
J.-W.~Chen, X.~Ji, and J.-H.~Zhang,
``Improved quasi parton distribution through Wilson line renormalization,''
Nucl.\ Phys.\ B \textbf{915}, 1 (2017), arXiv:1609.08102.
\quad\textbf{[来源:} \texttt{docs/Improved quasi parton distribution through Wilson line renormalization.pdf}\textbf{]}
（含辅助$z$-场形式——HQET启发的Wilson线重整化）

\bibitem{Zhang2019}
J.-H.~Zhang, X.~Ji, A.~Sch\"afer, W.~Wang, and S.~Zhao,
``Accessing gluon parton distributions in large momentum effective theory,''
Phys.\ Rev.\ Lett.\ \textbf{122}, 142001 (2019), arXiv:1808.10824.
\quad\textbf{[来源:} \texttt{docs/Accessing gluon parton distributions in large momentum effective theory.pdf}\textbf{]}
（含辅助伴随重夸克场方法——HQET向胶子算符的推广）

\bibitem{Zhang2023RGR}
R.~Zhang, J.~Holligan, X.~Ji, and Y.~Su,
``Leading power accuracy in lattice calculations of parton distributions,''
Phys.\ Lett.\ B \textbf{844}, 138081 (2023), arXiv:2305.05212.
\quad\textbf{[来源:} \texttt{docs/Leading Power Accuracy in Lattice Calculations of Parton Distributions.pdf}\textbf{]}
（含重夸克``极点''质量与重正子的类比讨论）

\bibitem{Chen2025gluon}
C.~Chen \textit{et al.} (CLQCD, Lattice Parton),
``Unpolarized gluon PDF of the nucleon from lattice QCD in the continuum limit,''
(2025), arXiv:2510.26425.
\quad\textbf{[来源:} \texttt{docs/Unpolarized gluon PDF of the nucleon from lattice QCD in the continuum limit.pdf}\textbf{]}

\bibitem{NoteGluonPDFs}
``Note of gluon PDFs,'' 内部笔记。
\quad\textbf{[来源:} \texttt{docs/Note of gluon PDFs.pdf}\textbf{]}

% ===== HQET经典文献 =====
\bibitem{ManoharWise}
A.~V.~Manohar and M.~B.~Wise,
\textit{Heavy Quark Physics}, Cambridge Monographs on Particle Physics, Nuclear Physics and Cosmology, Vol.~10 (Cambridge University Press, 2000).

\bibitem{Neubert1994}
M.~Neubert,
``Heavy quark symmetry,''
Phys.\ Rept.\ \textbf{245}, 259--396 (1994), arXiv:hep-ph/9306320.

\bibitem{IsgurWise}
N.~Isgur and M.~B.~Wise,
``Weak decays of heavy mesons in the static quark approximation,''
Phys.\ Lett.\ B \textbf{232}, 113--117 (1989).

\end{thebibliography}

\end{document}
